\chapter{Filters, Matching Networks, and Transformers}
\section{Why Filters Matter in Amateur Radio}
Filters select desired signals and reject unwanted ones. They appear in receiver
front ends, IF stages, audio stages, transmitter outputs, anti-aliasing paths,
power supplies, and interference cures.

\section{Basic Filter Types}
\begin{description}
\item[Low-pass] Passes low frequencies and attenuates high frequencies.
\item[High-pass] Passes high frequencies and attenuates low frequencies.
\item[Band-pass] Passes a frequency band.
\item[Band-stop or notch] Rejects a frequency band.
\end{description}

\begin{figure}[htbp]
\centering
\begin{circuitikz}
\draw (0,0) node[left]{in} to[R,l=\(R\)] (3,0)
      to[C,l=\(C\)] (3,-2) node[ground]{};
\draw (3,0) to[short,-o] (4,0) node[right]{out};
\end{circuitikz}
\caption{First-order RC low-pass filter.}
\end{figure}

\section{Cutoff and Slope}
For a first-order RC filter,
\[
f_c=\frac{1}{2\pi RC}.
\]
Past its corner, a first-order pole changes the asymptotic magnitude by about
\(-\SI{20}{\decibel}\) per decade. A second-order filter changes by about
\(-\SI{40}{\decibel}\) per decade after the relevant corner or resonant region.

\section{Filter Families}
\begin{tabularx}{\textwidth}{@{}L{0.2\textwidth}Y@{}}
\toprule
Family & Practical meaning \\
\midrule
Butterworth & Maximally flat passband; moderate transition. \\
Chebyshev & Ripple in the passband for a sharper cutoff. \\
Elliptic & Ripple in passband and stopband, with transmission zeros and a very sharp transition. \\
Crystal or cavity & High-Q resonators used when narrow bandwidth or strong selectivity is needed. \\
DSP FIR & Can be designed for linear phase; uses delayed samples and taps. \\
DSP IIR & Efficient recursive filter; phase is usually less linear than FIR. \\
\bottomrule
\end{tabularx}

\section{Impedance Matching}
An impedance-matching network does two jobs:
\begin{enumerate}
\item cancel the reactive part of a complex load, and
\item transform the remaining resistance to the desired system impedance.
\end{enumerate}
In a \(\SI{50}{\ohm}\) RF system, matching helps power transfer and keeps
transmission-line reflections controlled. Matching networks are usually narrow
band unless deliberately designed otherwise.

\section{Transformers}
For an ideal transformer,
\[
\frac{V_s}{V_p}=\frac{N_s}{N_p},
\qquad
\frac{I_s}{I_p}=\frac{N_p}{N_s}.
\]
The impedance reflected from secondary to primary is
\[
Z_p=\left(\frac{N_p}{N_s}\right)^2Z_s.
\]
Core saturation should be avoided because the magnetic core can no longer
support proportional flux change, causing distortion, heating, and loss of the
intended impedance transformation.

\section{Common Mistakes}
\begin{mistakebox}
Matching is not the same thing as filtering. A network can do both, but the
concepts are distinct: matching controls impedance seen by a source or line,
while filtering controls frequency-dependent transmission.
\end{mistakebox}

\section{Worked Example}
To match a real \(\SI{200}{\ohm}\) load to a \(\SI{50}{\ohm}\) system with a
lossless quarter-wave transformer,
\[
Z_t=\sqrt{Z_0Z_L}=\sqrt{(50)(200)}=\SI{100}{\ohm}.
\]
This formula assumes a real load and a transformer section that is one quarter
electrical wavelength long at the design frequency.
